\lysh{38875}{4}

\begin{alist}
\item
\begin{rlist}
\item Η παίκτρια, μέχρι την πέμπτη ερώτηση έχει $10 \cdot 3 - 2 \cdot 3 = 30 - 6 = 24$ βαθμούς.
Πετυχαίνει τη μέγιστη βαθμολογία, αν απαντήσει σωστά τις υπόλοιπες $25$ ερωτήσεις.
Αυτή είναι $24 + 25 \cdot 10 = 274$ βαθμούς.
\item Εφόσον δεν χάνει, απαντά σε όλες τις υπόλοιπες $25$ ερωτήσεις, είτε σωστά είτε
λάθος.
Η βαθμολογία που παίρνει από τις σωστές απαντήσεις είναι $10x$, ενώ από τις υπόλοιπες
$25 - x$ λανθασμένες χάνει $3(25 - x)$ βαθμούς.
Άρα η τελική βαθμολογία της είναι:
\[24 + 10x - 3(25 - x) = 24 + 10x - 75 + 3x = 13x - 51.\]
\item Εφόσον δεν έχασε και έφτασε μέχρι το τέλος του παιχνιδιού, η βαθμολογία της δεν
έγινε αρνητική κατά τη διάρκεια του παιχνιδιού.
Έτσι, μπορεί να χρησιμοποιηθεί ο τύπος του ερωτήματος ii. Άρα πρέπει:

$13x - 51 \ge 0$ ή $13x \ge 51$ ή $x \ge \dfrac{51}{13} \simeq 3,92$.

Όμως η τιμή του $x$ είναι ακέραιος αριθμός, γιατί εκφράζει πλήθος σωστών απαντήσεων.
Επομένως είναι $x \ge 4$, με $x$ ακέραιο.
Για $x = 4$ έχει τις λιγότερες σωστές απαντήσεις, άρα και την μικρότερη βαθμολογία, η
οποία είναι $13 \cdot 4 - 51 = 52 - 51 = 1$ βαθμός.
Πράγματι, αν απαντήσει σε $4 + 3 = 7$ ερωτήσεις σωστά και στις υπόλοιπες
λανθασμένα, τότε θα έχει $7 \cdot 10 - 23 \cdot 3 = 70 - 69 = 1$ βαθμό.
\end{rlist}

\item
\begin{rlist}
\item Με τη δεύτερη στρατηγική:
Εφόσον γνώριζε με βεβαιότητα $\sigma$ απαντήσεις, για αυτές πήρε $10\sigma$ βαθμούς.
Αν από τις υπόλοιπες $30 - \sigma$ ερωτήσεις έχει απαντήσει σωστά τις $x$, τότε:
\begin{itemize}
\item Ισχύει $0 \le x \le 30 - \sigma$.
\item Κέρδισε $10x$ βαθμούς, επιπλέον.
\item Απάντησε λάθος σε $30 - \sigma - x$ ερωτήσεις και για αυτές έχασε $3 \cdot (30 - \sigma - x)$
βαθμούς.
\item Για τις $30 - \sigma$ βαθμολογήθηκε με:
\[10x - 3(30 - \sigma - x) = 10x - 90 + 3\sigma + 3x = 13x + 3\sigma - 90.\]
\end{itemize}

Επομένως η συνολική του βαθμολογία για τις $30$ ερωτήσεις είναι:
\[10\sigma + 13x + 3\sigma - 90 = 13x + 13\sigma - 90.\]
\item Με την πρώτη στρατηγική θα έπαιρνε $4 \cdot 10 = 40$ βαθμούς.
Με τη δεύτερη στρατηγική, εφαρμόζοντας τον τύπο του $\beta$)i για $\sigma = 4$, βρίσκουμε ότι θα
έπαιρνε $13x + 13 \cdot 4 - 90 = 13x + 52 - 90 = 13x - 38$ βαθμούς, όπου $0 \le x \le 26$ με $x$ το πλήθος των
σωστών απαντήσεων, για τις οποίες δεν ήταν σίγουρος, όταν τις έδωσε.
Για να είναι επιτυχής η επιλογή στρατηγικής, πρέπει:
$13x - 38 \ge 40$ ή $13x \ge 78$ ή $x \ge 6$.
Άρα, αν πέτυχε επιπλέον $6$ ή περισσότερες σωστές απαντήσεις, εκτός από τις $4$ για τις
οποίες ήταν βέβαιος, η επιλογή στρατηγικής ήταν επιτυχής.
\end{rlist}
\end{alist}

